\clearpage
\VerdierPhysicalPageAnchor{0172}%
\begingroup\small\setlength{\parskip}{0pt}\setlength{\emergencystretch}{1.5em}
\VerdierSemanticAnchor{ch4-4-4}
\medskip\noindent\textbf{4.4. 收敛。}

\VerdierSemanticAnchor{ch4-4-4-1}
\medskip\noindent\textbf{4.4.1.}\quad
关于谱序列收敛性的详细研究，我们请读者参见\emph{同上引文}。这里只使用
如下判据：

\VerdierSemanticAnchor{ch4-4-4-2}
\medskip\noindent\textbf{定义 4.4.2.}\quad
设谱对象 $(H^n(p,q),\delta)$ 取值于阿贝尔范畴 $\mathcal A$。若对每个整数
$n$，都存在两个整数 $u(n),v(n)\in\mathbb Z$，使得：
\[
\begin{aligned}
H^n(p,q)&=0,&&u(n)\leq p\leq q\leq+\infty,\\
H^n(p,q)&=0,&&-\infty\leq p\leq q\leq v(n).
\end{aligned}
\]
则称该谱对象为\emph{定常的}。

\VerdierSemanticAnchor{ch4-4-4-3}
\medskip\noindent\textbf{定理 4.4.3.}\quad
\textit{设 $(H^n(p,q),\delta)$ 为定常谱对象。}

\textit{a) 对每个 $n$，滤过
$F^pH^n(-\infty,+\infty)$，$p\in\widetilde{\mathbb Z}$，都是定常的；
即存在两个整数 $n_1\leq n_2$，它们属于 $\mathbb Z$，使得：}
\[
F^{n_1}H^n(-\infty,+\infty)=0
\quad\text{且}\quad
F^{n_2}H^n(-\infty,+\infty)=H^n(-\infty,+\infty).
\]

\textit{b) 从
$E\!\left[\begin{smallmatrix}n\\r,p,q,s\end{smallmatrix}\right]$ 出发
（相应地，以其为靶）的微分，在 $r\leq v(n+1)$（相应地，
$u(n-1)\leq s$）时为零；而且：}

\textit{c) 对每个 $n$ 及每个数对 $p\leq q$，对象
$E\!\left[\begin{smallmatrix}n\\r,p,q,s\end{smallmatrix}\right]$ 若满足}
\[
-\infty\leq r\leq v(n+1)
\quad\text{且}\quad
u(n-1)\leq s\leq+\infty
\]
\textit{则彼此典范同构。}

本定理由定义立即推出。

\VerdierSemanticAnchor{ch4-4-4-3-convergence}
一个定常谱对象给出一个在\emph{同上引文}意义下强收敛的谱序列。当与谱对象
关联的谱序列所对应的谱对象定常时，我们称该谱序列为\emph{定常的}，或宽泛
地称为\emph{收敛的}。采用 (4.3.3) 中的任一种记号，我们写成：
\[
\mathrm I_r^{p,q}\Longrightarrow H^{p+q}(-\infty,+\infty),
\qquad
\mathrm{II}_r^{p,q}\Longrightarrow H^{p+q}(-\infty,+\infty)
\]
以表示一个谱序列收敛。此时，各项 $H^{p+q}(-\infty,+\infty)$ 的全体
称为该谱序列的\emph{收敛目标}。
\endgroup
